% ============================================================
%  Chapter 5 — Evaluation and Results
% ============================================================
\chapter{Evaluation and Results}
\label{chap:evaluation}

\section{Experimental Setup}
\label{sec:setup}

To evaluate the effectiveness of structural context retrieval, we use the \textbf{SWE-bench}~\cite{yang2024swebench} framework, which provides a realistic testbed for software engineering tasks.

\subsection{Benchmark: SWE-bench Lite}
SWE-bench Lite is a curated subset of 300 instances drawn from real-world Python repositories (e.g., Django, SymPy, Matplotlib). Each instance consists of:
\begin{itemize}
    \item \textbf{An Issue Description:} A natural language bug report or feature request.
    \item \textbf{A Repository Snapshot:} The state of the codebase before the patch.
    \item \textbf{A Ground Truth Patch:} The files that must be modified to resolve the issue.
\end{itemize}

\subsection{Metrics}
The primary objective of our retrieval system is to ensure the ground-truth files are present in the context provided to the agent. We measure performance using:
\begin{itemize}
    \item \textbf{Recall@$k$:} The percentage of instances where at least one ground-truth file is present in the top-$k$ retrieved results.
    \item \textbf{Mean Reciprocal Rank (MRR):} A measure of how high in the list the first relevant file appears.
\end{itemize}

\subsection{Baseline: BM25}
We compare our approach against a pure keyword-based baseline using BM25. This baseline represents the current industry standard for RAG-based coding assistants that do not leverage repository-wide structural information.

\section{Results}
\label{sec:results}
% Present key findings from benchmark iterations.

\subsection{Iteration 1: Weighted PPR}
% The initial large gain from implementing weighted PPR.

\subsection{Iteration 2: Refined Seed Selection}
% Further improvement from refining seed selection.

\subsection{Iteration 3: Structural Limitations}
% The finding that adding more structural edges (CO_LOCATED) did not help, and why.

\section{Analysis of Results}
\label{sec:analysis}
% Interpret the data. Why did PPR outperform the baseline? What does Iteration 3 imply?
